\documentclass[aspectratio=169,11pt]{beamer}
\usepackage[utf8]{inputenc}
\usepackage[T1]{fontenc}
\usepackage[french]{babel}
\usepackage{helvet}
\renewcommand{\familydefault}{\sfdefault}
\usepackage{hyperref}
\definecolor{bleu}{HTML}{0000FF}
\definecolor{fond}{HTML}{F4F4EF}
\setbeamercolor{background canvas}{bg=fond}
\setbeamercolor{structure}{fg=bleu}
\setbeamercolor{title}{fg=bleu}
\setbeamercolor{frametitle}{fg=bleu,bg=fond}
\setbeamercolor{normal text}{fg=black}
\setbeamertemplate{navigation symbols}{}
\setbeamertemplate{itemize item}{\textcolor{bleu}{\tiny$\blacksquare$}}
\setbeamertemplate{footline}{\hfill\footnotesize\insertframenumber\hspace{1em}\vskip6pt}
\hypersetup{colorlinks=true,urlcolor=bleu,linkcolor=bleu}
\title{Intégration économique}
\subtitle{Union africaine}
\author{Veille institutionnelle}
\date{2026-11-03}

\begin{document}
\begin{frame}\titlepage\end{frame}

\begin{frame}{Intégration économique}
\begin{itemize}
\item ZLECAf : zone de libre-échange continentale, en vigueur en 2019, échanges depuis 2021, secrétariat à Accra
\item Communautés économiques régionales reconnues par l'UA : huit (à vérifier)
\item AUDA-NEPAD : agence de développement de l'UA
\item Protocole sur la libre circulation des personnes
\end{itemize}
\end{frame}

\begin{frame}{Intégration économique : ressources à suivre}
\small
Veille hebdomadaire.
\begin{itemize}
\item \href{https://au-afcfta.org}{ZLECAf}\\{\footnotesize Site spécialisé, publication : hebdomadaire}
\item \href{https://au.int/en/recs}{Communautés économiques régionales}\\{\footnotesize Page institutionnelle, publication : trimestrielle}
\item \href{https://www.nepad.org}{AUDA-NEPAD}\\{\footnotesize Site spécialisé, publication : mensuelle}
\end{itemize}
\end{frame}

\end{document}
